\section{Acquiring a homogeneous phase modulo quotient noise}\label{sec:quotientfit}
This section proves the acquisition part of Theorem~\ref{thm:quotient}.
The full third tensor of \eqref{eq:quotient-family} is the scalar volume
form with unit coefficient $\mu$. Repeated-variable ordinary cubic
monomials, including $X^3$, have binary degree at most two in
characteristic two. Thus all third coefficients determine exactly the
term $\lambda(\mu XYZ)$, while the remaining Boolean quadratic is explicit.
The preceding centroid construction therefore acquires $A$ and its
rank-three free module without a supplied ring presentation.

\subsection{Why the promise survives acquired coordinates}
Two presentations of this full centroid differ by a ring isomorphism
$\sigma$ and an invertible module matrix. Including translations, the
change has the form $x=M\sigma(y)+b$, with $M\in\GL_3(A)$.
Nondegeneracy of the Frobenius pairing gives a unique unit $c$ such that
the transported functional is $z\mapsto\lambda'(cz)$. Consequently
the transformed ordinary homogeneous cubic is
$c\sigma^{-1}(H\circ M)$, with a translated center.
Its mixed coefficient is
$c\sigma^{-1}(\mu\det M)$, again a unit. A repeated-variable cubic
cannot create an $XYZ$ term under linear substitution in characteristic two.
The ideal $J^h$ is preserved by both $M$ and $\sigma$.
The transported $Q$ is still a binary quadratic on $A^3/(J^h)^3$.
This argument includes nonlinear changes of uniformizer and coefficient-field
Frobenius; no preferred hidden basis is assumed.

\begin{lemma}[Ordinary coefficients invisible above a quotient]\label{lem:invisible}
Let $d\ge3$. Suppose the Boolean function $\lambda(P(x))$, where $P$ is
ordinary of total degree at most three, depends only on $x\bmod J^h$.
After replacing ordinary square terms by their exact linear adjoints,
every homogeneous cubic coefficient and every cross-quadratic
coefficient of $P$ lies in $J^{e-h}$.
\end{lemma}
\begin{proof}
For all $x$ and $v\in(J^h)^3$, the difference
$\lambda(P(x+v)+P(x))$ is zero. Scale both $x$ and $v$ by every
$\rho\in F$. Writing the homogeneous pieces of $P$ separately gives
\[
 \Tr\bigl(a\rho^3+b\rho^2+c_1\rho\bigr)=0
 \quad\text{for every }\rho\in F,
\]
where $a=[t^{e-1}](P_3(x+v)+P_3(x))$.
For $d\ge3$, the doubling orbit of exponent $3$ modulo $2^d-1$ has
length $d$ and is disjoint from the orbit of $1$ (which contains $2$).
Hence the trace polynomial identity forces $a=0$.
This is an ordinary polynomial identity over $F$ in the coefficient
variables of $x$ and $v$; its degree is at most three and $|F|\ge8$,
so its coefficients vanish.

Each degree-three monomial has an odd exponent. To isolate its coefficient
at $t^j$, for $j<e-h$, use once the variable $v_i$ at power
$t^{e-1-j}$ in such an odd-exponent slot, and use residue variables
for the other occurrences. The binomial coefficient is odd and no other
monomial contributes the same coefficient-variable monomial. Thus that
coefficient is zero. This proves $P_3\in J^{e-h}A[X,Y,Z]$.
It already defines a function on the low quotient, so subtract it.
Binary polarization between different module coordinates and
nondegeneracy of the field trace pairing then force every cross-pair
coefficient into $J^{e-h}$ as well. Square terms have been folded into
linear terms and do not enter this last polarization. Once the cubic and
cross terms are removed, the same Frobenius pairing forces every remaining
linear coefficient into $J^{e-h}$ as well.
\end{proof}

\subsection{A linear fit that does not guess the error}
Set $h=\lfloor(e+1)/3\rfloor$; this includes every smaller allowed quotient.
Keep the acquired term $\mu XYZ$ fixed. For the other nine ordinary
homogeneous cubic monomials, three cross pairs, three linears and a
constant, form their binary trace-ANF columns, one for each ring
coefficient bit. There are $16ed$ columns. All have binary degree at
most two, so every acquired third coefficient remains fixed.
Constrain every quadratic or linear ANF coefficient involving a bit
outside the low quotient. The remaining rows may be arbitrary noise.
The protected coefficient matrix has exact rank $15d(e-h)$: its kernel
consists of arbitrary constant coefficients and the $15$ nonconstant
coefficient families in $J^{e-h}$, hence has dimension $de+15dh$.
Solve this binary linear system and obtain any ordinary polynomial $P$.
Failure rejects the input.

Let $H=P_3$ and read the affine center directly:
\[
 a=(P_{YZ},P_{XZ},P_{XY})/\mu.
\]
Subtract $\lambda(H(x+a))$ from the original acquired phase. Its third
tensor is zero, so all remaining coefficients are recovered exactly
from evaluations of weight at most two. Verify that every resulting
monomial is supported on the low quotient. Finally translate $x$ by
$a$ and carry the translated quotient quadratic $Q'$ explicitly.
The output is the exact identity
\[
 p(x+a)=\lambda(H(x))+Q'(x\bmod J^h).
\]

To prove completeness, compare any fitted $P$ with the ordinary trace
polynomial obtained from the promised translated homogeneous phase.
Their trace difference depends only on the quotient. By
Lemma~\ref{lem:invisible}, their homogeneous cubics and cross-pair
coefficients differ by elements of $J^{e-h}$. Their centers therefore
differ by $\delta\in(J^{e-h})^3$. Substituting such a shift in a cubic
changes its trace only through the low quotient: coefficients of every
new nonconstant term contain a factor from $J^{e-h}$, while terms
with two such factors vanish because $e\ge2h$.
Thus every linear-system solution produces an admissible normalized
identity. This removes a search among fit solutions.
The final coefficient check also certifies accepted inputs outside the
promised generation history.

\section{Exact counting by manufactured critical rays}\label{sec:criticalrays}
Now $H$ and $Q$ are acquired explicitly. No factorization of $H$ is
needed. Write $q=2^d$, $\psi_e(z)=(-1)^{\lambda(z)}$, and let
$H_0$ denote reduction modulo $t$.
For $e\ge2$, averaging over the highest digit of each variable cancels
every nonzero residue vector where $\nabla H_0\ne0$. The quotient
quadratic is independent of those highest digits. The surviving
nonzero residue classes lie on projective singular rays of $H_0=0$.
Euler's degree-three identity in characteristic two gives
$H_0(x)=x\cdot\nabla H_0(x)$, so a gradient zero lies on the curve.

\subsection{Finding every ray by binary linear algebra}
Use disjoint projective charts: $Z\ne0$ first; then $Z=0,Y\ne0$;
then $Z=Y=0,X\ne0$. In the first chart put $Z=1$ and
$f(u,v)=H(u,v,1)$. Its gradients have the form
\[
 f_u=a u^2+b v^2+\mu v+c_1,\qquad
 f_v=c u^2+d_1 v^2+\mu u+c_2.
\]
Over the residue field this is an affine binary linear system with
$2d$ unknown bits. Its Jacobian is
$\left(\begin{smallmatrix}0&\mu_0\\\mu_0&0\end{smallmatrix}\right)$,
with nonzero determinant. Two quadratics with this Jacobian have at
most four geometric solutions; hence the binary solution space has
kernel dimension at most two. Enumerate these at most four candidates,
apply the chart restrictions, and retain $H_0=0$.

There are at most three projective singular rays. A repeated plane factor
would make $H_0=L^2M$ and annihilate its $XYZ$ coefficient. Thus the
cubic is reduced. A line through two singular points divides a cubic,
since the two intersection multiplicities already total at least four.
Three distinct singular points therefore force a triangle of lines;
there cannot be a fourth. This also covers nonsplit residue curves.

Each transverse critical point has a unique lift $a\in A^2$, because
the displayed Jacobian is a unit. At layer $t^k$ solve the two residual
gradient coefficients using its inverse. All operations are in the
acquired ring. The lifted vector $v=(a_1,a_2,1)$, the critical value
$c=f(a)$, and the residue quadratic coefficients below are manufactured.
Their truncations are the unique lifts at every smaller depth, so they
can be reused throughout the origin recursion.

\subsection{The exact transverse calculation in characteristic two}
Translate to the lifted critical point. The chart phase becomes
\[
 f(a+w)=c+A w_1^2+\mu w_1w_2+B w_2^2+C_3(w_1,w_2),
 \qquad w\in(tA)^2,
\]
where $C_3$ is homogeneous cubic. Let $s=\lfloor e/2\rfloor$ and
$k=\lceil e/2\rceil$. Partition the transverse domain into cosets of
$(t^kA)^2$. Addition of $\delta\in(t^kA)^2$ changes the phase exactly
by its formal gradient paired with $\delta$, because $\delta^2=0$ in
$A$. The weight $Q$ is unchanged, since $h\le s\le k$.
Character orthogonality therefore forces
$\nabla f(a+w)\in(t^sA)^2$.

The gradient is an invertible linear term $\mu(w_2,w_1)$ plus a
quadratic-in-$w$ term. On $(tA)^2$ it preserves the minimum valuation:
for a nonzero vector of valuation $j\ge1$, the linear term has
valuation $j$ and the quadratic term at least $2j$.
Thus the forced condition is exactly $w\in(t^sA)^2$.
For even $e=2s$, all nonconstant terms then vanish and the transverse
sum is $q^e$. For odd $e=2s+1$, the highest free digits contribute
$q^{2s}$ and the two middle digits contribute the field Gaussian
\[
 \sum_{u,v\in F}(-1)^{\Tr(A_0u^2+\mu_0uv+B_0v^2)}
   =q\kappa,\qquad
 \kappa=(-1)^{\Tr(A_0B_0/\mu_0^2)}.
\]
The cubic term vanishes because $3s\ge e$ for odd $e\ge3$.
For a fixed radial unit $r$, every coefficient is multiplied by $r^3$,
which cancels from the displayed Arf ratio. The quotient weight is
constant on this surviving transverse cell and equals $(-1)^{Q(rv)}$.
The exact contribution of one projective ray is consequently
\begin{equation}\label{eq:weightedray}
 q^e\kappa^{e\bmod2}
 \sum_{r\in A^\times}(-1)^{\lambda(c r^3)+Q(rv)}.
\end{equation}
No characteristic-zero Morse lemma is used.

\subsection{A quadratic radial phase and a single recursive state}
The map $r\mapsto r^2$ is binary linear. Hence
$\lambda(c r^3)+Q(rv)$ is a binary quadratic on $ed$ bits, even when
$Q$ has large rank. The unit sum in \eqref{eq:weightedray} is exactly
\[
 \mathcal G_A\bigl(\lambda(c r^3)+Q(rv)\bigr)
 -\mathcal G_{tA}\bigl(\lambda(c r^3)+Q(rv)\bigr),
\]
where each $\mathcal G$ is computed by ordinary binary symplectic
elimination, including its radical and affine term. The two Gaussian
variable counts are $ed$ and $(e-1)d$. No character expansion of $Q$
is formed. For $Q=0$, the unit sum also has the direct formula
\[
 U_e(c)=\begin{cases}
 0,&v_t(c)\le e-2,\\
 q^{e-1}\left(\sum_{r\in F}(-1)^{\Tr(br^3)}-1\right),
       &c=t^{e-1}b.
 \end{cases}
\]
This formula includes $c=0$. The field cube phase is binary quadratic;
no claim that cubing permutes $F^\times$ is needed.

The only remaining residue class is the origin. Put $x=ty$.
For $e\ge4$, homogeneity gives the exact origin contribution
\[
 q^6 Z_{e-3}\bigl(H\bmod t^{e-3},\ Q(ty)\bigr).
\]
There are two unused ring digits in each of three variables. The
quotient height decreases by one, and
$3(h-1)\le(e-3)+1$, so the same hypothesis persists.
For $e=2$ or $3$, the homogeneous cubic vanishes on the origin class
and the contribution is
$(-1)^{Q(0)}q^{3(e-1)}$.
At most three nonrecursive ray contributions are added at each level.
Only one origin state recurses, for at most $\lceil e/3\rceil$ levels.
This is why the number of computations is polynomial rather than an
exponentially branching lift tree.

\subsection{The finite-field terminal}\label{sec:fieldterminal}
When the origin recursion reaches depth one, the allowed quotient
height is zero. The remaining phase is
$\Tr(H_0(X,Y,Z))$ plus a constant. Manufacture the ordinary equation
\[
 V:\quad W^2+W+H_0(X,Y,Z)=0
\]
over $F$. Each triple has two lifts $W$ when its trace is zero and
none otherwise, so
\[
 Z_1=\#V(F)-q^3.
\]
This is a four-variable, degree-three, characteristic-two point count.
Lauder--Wan, Proposition~36 and the affine torus decomposition in
Section~6.4, give deterministic bit time $\widetilde O(d^{19})$ and
space $\widetilde O(d^{12})$ for this first point count \cite{lauderwan}.
The constants depend only on the fixed degree, dimension and characteristic;
no smoothness is required. Computing an entire zeta function is unnecessary.

The generic point-count engine is not implemented in this source package.
The executed code explicitly limits its field terminal to $3d\le24$
bits by default, and records every enumerated point. For
$e\equiv0,2\pmod3$, it never reaches a field terminal and has no
field-size restriction. The separate file-leaf compiler manufactures
and verifies the full trace-polynomial identity using $20d$ binary
coefficient columns. It rejects arbitrary binary quadratics that do
not have an ordinary degree-three field lift. Its Artin--Schreier
reduction is executed independently; it is not advertised as an
execution of the general Lauder--Wan algorithm.

An additional implemented specialization handles the field family
$X^2Y+XY^2+Z^3+aXYZ+\rho X+\sigma Y+\tau Z+c$.
Eliminating $X$ and substituting $Y=w^2$ leaves a Laurent trace phase
with exponents among $\{3,1,-1,-3,-5\}$, plus an exact Gaussian
boundary contribution. The finite Dwork trace backend has at most
$32(d+2)-1$ states and conservative $O(d^8)$ bit time and $O(d^4)$
space. This is an executed specialization of established Dwork machinery.
No acquired equivalence from every cubic to that special form is claimed.

\subsection{Description size and arithmetic costs}
The acquired binary coordinate map uses $n^2$ bits. Ten ring coefficients
use $10ed$ bits. A quotient quadratic uses at most
$1+3dh+\binom{3dh}{2}$ binary coefficients, or $O(n^2\log(n+2))$
bits in an index-list encoding. The complete representation is therefore
polynomial in the original independent-input count. It is manufactured
from all raw coefficients, not passed to the counter as advice.

For an explicit conservative cost, binary centroid elimination has
$n^2$ unknowns and $O(n^3)$ equations, and costs at most $O(n^7)$
bit operations with classical elimination. Ring structure, Frobenius
subspaces and free module bases use polynomial-size binary matrices;
the algorithms described above fit within $O(n^8)$ bit time.
The quotient fit has $O(n)$ columns and $O(n^2)$ rows. Each candidate
trace-ANF is obtained from ring monomial evaluations and a Gram matrix;
all ring elements have at most $n$ bits. Its construction, solve,
normalization and exact coefficient checks fit within the same bound.

There are $O(e)$ stages and at most six binary Gaussian sums per stage.
Each has at most $ed\le n$ variables. Even constructing its coefficients
by all weight-at-most-two evaluations, each evaluation scanning an
$O(n^2)$-term quadratic with $n$-bit masks, gives a bound below $O(n^8)$
for the entire recurrence. Schoolbook multiplication of ring coefficients
and of count integers is included. Streaming the binary equations and
retaining polynomial-size matrices gives $O(L+n^4)$ space before a
finite-field terminal. Combining these bounds with the displayed
Lauder--Wan term proves Theorem~\ref{thm:quotient}.

All count contributions and recursion multipliers have $O(n)$ bits.
The original signed sum lies between $-2^n$ and $2^n$; the final zero
count is reconstructed as $(2^n+Z)/2$. No unit-cost arithmetic over
exponentially large integers is assumed. For an accepted gate program,
add the explicit gate-to-ANF cost from the preceding section; internal
wire values are never counted as independent inputs.

\section{Smooth shortcuts, sharp failures, and an adaptive repair}
For a smooth residue cubic there are no nonzero critical rays. The
weighted recurrence therefore becomes exact digit forcing toward a
single translated origin. For a homogeneous centered phase and the
allowed quotient height,
\[
 Z_e(H,Q)=(-1)^{Q(0)}
 \begin{cases}
 q^{2e},&e\equiv0\pmod3,\\
 q^{2e-1},&e\equiv2\pmod3,\\
 q^{2(e-1)}Z_1(H_0),&e\equiv1\pmod3.
 \end{cases}
\]
The earlier smooth-only compiler manufactures the residue center directly
from highest-digit derivatives. After fitting each binary-linear map as
$M_{ij}X_j^2+B_{ij}X_j$, it reads the center
\[a=(B_{23},B_{13},B_{12})/\mu_0,\]
checks every gradient identity, and
certifies absence of nonzero homogeneous gradient zeros with an
$18\times15$ Macaulay matrix in degree four. Full rank means that the
three quadrics generate every quartic monomial. Conversely a zero-free
homogeneous gradient is a regular sequence of quadrics, whose quotient
has Hilbert series $(1+t)^3$, so this certificate is complete.
For $d=2$ the square and linear maps are still distinguishable; for $d=1$
the implementation checks the exact eight residue vectors instead.

The smooth-only implementation also checks that the proposed discarded
layers are absent. This is essential: over $\F_4[t]/t^4$ with
$H=X^2Z+Y^3+Z^3+XYZ$, adding one linear bit from the $t^2$ coefficient
leaves the first center and smoothness test unchanged but changes the
signed count from $114688$ to $0$. The unused-layer test rejects the
invalid compression. The independent audit executed both full truth tables.

The sign-only invisibility identity has a sharp height limit. Over
$\F_2[t]/t^e$ use the same smooth $H$ and set
$h_0=\lfloor(e+1)/3\rfloor$. A rank-two quadratic using the next digit
$t^{h_0}$ changes $Z$ by a factor $1/2$, $-1$, or $2$ according as
$e\equiv0,1,2\pmod3$: use $X_{h_0}Y_{h_0}$ in the first and third
cases and $X_{h_0}Z_{h_0}$ in the second. The exact terminal signed
sums at depths $3,1,2$ are respectively $64,2,8$, and become
$32,-2,16$. Valid origin reductions propagate these examples to all
depths. This refutes the tempting simplification, not polynomial counting
outside the sufficient height range.

Indeed the critical-ray calculation only needs $2h\le e$ at the current
stage. Continue it while this holds, then keep the single remaining
state explicitly. If
$\Delta=h-\lfloor(e+1)/3\rfloor>0$, the residual number of binary
variables is at most
\[
 R\le\min\{n,\ 3d\max(1,6\Delta-1)\}.
\]
At level $j$ the depth is $e-3j$ and the nonzero height is $h-j$;
the condition is $j\le e-2h$. This gives the displayed bound on the
first unhandled state. The executable accepts a separate residual-bit
budget and records its $2^R$ enumeration explicitly. The additional
cost is $2^R\poly(L+n)$, not exponential in the full rank of $Q$.
For fixed $d$ and fixed excess height, it is polynomial in growing $e$.

The raw fitter can also acquire the minimum admissible height. The
coefficient-ambiguity lemma holds for every $0\le h\le e$. Its center
argument does not actually need $\delta^2=0$: every changed coefficient
contains $\delta\in J^{e-h}$ and hence remains supported on the low
quotient, even when higher powers of $\delta$ survive.
Existence of the fitted representation is monotone in $h$, and height
$e$ always permits $H=\mu XYZ$ plus the full explicit quadratic.
Binary search using the exact fit therefore manufactures the smallest
height. The semilinear invariance argument makes this minimum independent
of the acquired chain presentation. No favorable height is
supplied. Combining it with the residual bound gives a certified
adaptive algorithm on every raw rank-three chain-volume cubic with
$d\ge3$, with graceful but explicitly exponential cost in an unfavorable
residual. It does not turn arbitrary dense quadratic couplings into a
polynomial family.

\section{Obstructions that the construction does not erase}
\subsection{Arbitrary linear constraints recover general cubic counting}\label{sec:parityboundary}
The unrestricted character sum of a ring or field phase can be easy while
its restriction to an arbitrary binary subspace is hard. This matters for
conditional queries, inversion, synthesis and autoregressive learning:
none follows from a total count without checking the new constraints.

Here is an explicit reduction. For any Boolean cubic $p$ on $m$ bits,
choose $F=\F_{2^{m^2}}$ with polynomial basis $1,\alpha,\ldots$.
Put $a_i=\alpha^i$ and $b_j=\alpha^{mj}$ for $0\le i,j<m$.
Their $m^2$ products are a basis. For each $k$, solve the nondegenerate
trace-pairing system for $c_k$ so that
$\Tr(a_i b_j c_k)$ is an assigned coefficient array. Assign each Boolean
monomial of $p$ to one ordered triple, permitting repeated indices for
linear and quadratic terms. Then the injective binary map
\[
 u\longmapsto\left(\sum_i a_i u_i,\sum_j b_j u_j,\sum_k c_k u_k\right)
\]
restricts $\Tr(XYZ)$ to $p$ (a separate constant sign is harmless).
A parity-check description of its image is manufactured by elimination
and uses $O(m^4)$ bits. The fixed-characteristic field construction is
polynomial; standard deterministic constructions are available \cite{shoup}.

Cubic signed counting is already \#P-hard: encode every Boolean circuit
gate by its unique-extension equation $f_j=0$ of degree at most two,
including an accepting-output equation, and introduce independent bits
$z_j$. Orthogonality gives
\[
 \sum_{x,\text{gate states},z}(-1)^{\sum_j z_j f_j}
   =2^{\#\{j\}}\#\{\text{accepting original inputs}\}.
\]
Thus a general polynomial counter for arbitrary parity-constrained
$\Tr(XYZ)$ would imply $\mathrm{FP}=\#\mathrm P$.
This is a reduction, not an assumed separation.

A separate Fourier calculation shows that a rank-$r$ quadratic sign has
exactly $2^r$ nonzero Walsh coefficients. Any universal reconstruction
using fixed linear-character shifts and weights needs that many terms.
This is a restriction on that representation, not a lower bound on
counting time. The weighted critical-ray method explicitly escapes it
for compatible quotient weights by using a different reconstruction.

\subsection{A nonchain centroid can retain an arbitrary tensor problem}
For any symmetric trilinear form $S$ on $F^k$, construct the commutative
Frobenius algebra
\[
 A_S=F\oplus U\oplus U^*\oplus F\omega,\qquad U=F^k.
\]
Define $uv=S(u,v,\cdot)\in U^*$, $u\phi=\phi(u)\omega$, and all other
positive-degree products by the grading. Symmetry of $S$ gives
associativity, the radical has fourth power zero, and the functional
extracting $\omega$ has a nondegenerate multiplication pairing.
The centroid of the scalar volume tensor on $A_S^3$ is exactly $A_S$.
Nevertheless its cubic signed sum satisfies the explicit identity
\[
 Z_{A_S}=q^{3k+3} Z_S+3(q-1)q^{4k+3},
\]
where $Z_S=\sum_{u,v,w\in F^k}(-1)^{\Tr S(u,v,w)}$.
Thus acquisition of an arbitrary local algebra can leave the original
trilinear counting problem inside it. The polynomial chain-ring theorem
cannot be extended merely by replacing the word ``chain'' with ``local''.
The algebra construction and the displayed identity have complete
proofs and executed multiplication/count checks in the source files.

\subsection{Other preserved failures}
A residue cubic may split into three independent planes without its full
ring lift splitting: the pure-cube coefficients in the Hensel residual
are genuine obstruction checks. The global nilpotent annihilator may
annihilate a defect in its kernel even though the residual is not an
ordinary ring cubic; the final coefficient fit is essential.
Small residue fields can identify ordinary cubic trace coefficients,
which is why the quotient-fit completeness argument retains $d\ge3$.
The square adjoint in \eqref{eq:square} must include nilpotent indices;
a field-only square-root formula is false.

Finally, the new quotient family permits arbitrary-rank compatible
quadratics, but not arbitrary ordinary linear perturbations. For example,
$\lambda(XYZ+X)$ at depth two fails the translated-homogeneous low-quotient
fit. It remains supported by the complementary split-cubic counter.
This failure is retained as an actual rejection example, not bypassed
by dropping unmatched coefficients.
